% Figure nearterm: generated by figures/rse_fig_nearterm.py; placed where the text first discusses it.
% Display block, shading construction, clipping and scale bar: Supplementary section rse_si_display.tex.
% The first-metre range of the five grids is given in the Results text (sec:res_nearterm).
\begin{figure}[!tbp]
  \centering
  \includegraphics[width=\textwidth]{../figures/out/rse_fig_nearterm.pdf}
  \caption{\textbf{The next metres, measured from the observed 2026 shoreline.}
  Each grid is offset so that its sea at the 2026 level (\metric{h3.level_2026_m}~m, DAHITI reference) has the area of the observed 2026 sea (\metricsf{rev2.n.reref.S_obs_km2}{3}~km\textsuperscript{2}); the observed sea that a grid still places on dry ground (the misfit) is \metricsf{rev2.n.reref.area.main.M.min_km2}{2}--\metricsf{rev2.n.reref.area.main.M.max_km2}{2}~km\textsuperscript{2} (\metricdp{rev2.n.reref.area.main.M.min_pct_of_S_obs}{1}--\metricdp{rev2.n.reref.area.main.M.max_pct_of_S_obs}{1}\%).
  (a)~Number of grids that place each cell of the observed 2026 sea in the first metre of further fall (uncoloured, none).
  At least one grid places \metricsf{rev2.n.agree60.observed_reref.1.any_km2}{3}~km\textsuperscript{2} there and all five agree on \metricsf{rev2.n.agree60.observed_reref.1.all_km2}{2}~km\textsuperscript{2} (\metricdp{rev2.n.agree60.observed_reref.1.all_pct_of_any}{1}\%).
  Light blue: the open sea of \metric{params.lowstand_label} from Landsat.
  (b)~Emerging seabed $E(D)$, the area that emerges between the 2026 shoreline and a further fall $D$, for each grid: solid up to \metric{rev.n.env.max_pair_decline_m}~m, the largest fall in the hindcast, dotted beyond (not calibrated).
  Shading: the range of $E(D)$ when the 2026 level is raised or lowered by \metric{rev.w.n.sens.jitter_m}~m.
  (c)~Fall below the 2026 level from 2026 to 2070 for the seven primary published pathways (each pathway's own change after 2026); grey band, the range pathways (inter-model ranges and quartiles); dashed line, the calibrated limit of b.}
  \label{fig:rse_nearterm}
\end{figure}
